\documentclass[11pt]{article}

% ── Packages ──────────────────────────────────────────────────────────
\usepackage[margin=1in]{geometry}
\usepackage{amsmath,amssymb,amsthm}
\usepackage{booktabs}
\usepackage{graphicx}
\usepackage{hyperref}
\usepackage{algorithm}
\usepackage{algorithmic}
\usepackage{multirow}
\usepackage{xcolor}
\usepackage{subcaption}
\usepackage{pgfplots}
\pgfplotsset{compat=1.18}

% ── Theorem environments ─────────────────────────────────────────────
\newtheorem{theorem}{Theorem}
\newtheorem{proposition}[theorem]{Proposition}
\newtheorem{lemma}[theorem]{Lemma}
\newtheorem{corollary}[theorem]{Corollary}
\newtheorem{definition}[theorem]{Definition}
\newtheorem{remark}[theorem]{Remark}

% ── Macros ────────────────────────────────────────────────────────────
\newcommand{\R}{\mathbb{R}}
\newcommand{\Z}{\mathbb{Z}}
\newcommand{\argmin}{\mathop{\mathrm{argmin}}}
\newcommand{\conv}{\mathrm{conv}}
\newcommand{\opt}{\mathrm{opt}}
\newcommand{\BiCut}{\textsc{BiCut}}

\title{\BiCut: A Solver-Agnostic Bilevel Optimization Compiler\\with Intersection Cuts for MIBLPs}

\author{BiCut Contributors}

\date{}

\begin{document}
\maketitle

\begin{abstract}
We present \BiCut, a typed bilevel optimization compiler that transforms
mixed-integer bilevel linear programs (MIBLPs) into solver-ready
single-level reformulations with machine-checkable correctness certificates.
\BiCut\ introduces \emph{bilevel intersection cuts}---a new family of valid
inequalities derived by extending Balas's (1971) intersection cut framework
to the bilevel-infeasible set defined by follower suboptimality.
For LP lower levels, we prove that the bilevel-infeasible set is a finite
union of polyhedra, establish a polynomial-time separation oracle (for fixed
follower dimension), and prove finite convergence of the cut loop under
non-degeneracy. The compiler supports four reformulation strategies
(KKT, strong duality, value function, column-and-constraint generation),
three complementarity encodings, and emits to Gurobi, SCIP, and HiGHS.
On our current generated suite of 122 bilevel instances, the KKT
reformulation solved by HiGHS performs well (121/122 instances solved,
most at the root node). The intersection cut framework is
theoretically grounded, but in this preliminary study the cuts do not
yet improve upon HiGHS's native cut machinery on these instances.
\BiCut\ is the first bilevel compiler with correctness certificates and
the first tool providing bilevel-specific cutting planes.
\end{abstract}

% ══════════════════════════════════════════════════════════════════════
\section{Introduction}
\label{sec:intro}

Bilevel optimization---where one optimization problem is nested as a
constraint within another---arises in network interdiction
\cite{israeli2002shortest}, competitive facility location
\cite{kucukaydin2011competitive}, strategic bidding in electricity markets
\cite{ruiz2009pool}, Stackelberg security games
\cite{tambe2011security}, and adversarial robustness certification
\cite{sinha2018review}. The canonical mixed-integer bilevel linear
program (MIBLP) has the form:
\begin{align}
\min_{x \in X} \quad & d^\top x + e^\top y \label{eq:bilevel}\\
\text{s.t.} \quad & Cx + Dy \leq h, \nonumber\\
& y \in S(x) = \argmin_{y \geq 0}\{ c^\top y : Ay \leq b + Bx \}, \nonumber
\end{align}
where the \emph{leader} controls $x$ and the \emph{follower} responds
optimally with $y = S(x)$.

Despite two decades of algorithmic progress on MIBLPs
\cite{dempe2002foundations, colson2007overview}, the software landscape
remains fragmented. Existing tools---BilevelJuMP.jl \cite{diasgarcia2022bilevel},
PAO \cite{hart2017pyomo}, GAMS EMP, YALMIP \cite{lofberg2004yalmip}---support
one or two reformulation strategies with no automatic selection, no correctness
guarantees, and no principled handling of mixed-integer lower levels.

This paper makes two contributions:

\begin{enumerate}
\item \textbf{Bilevel intersection cuts.} We extend Balas's 1971 intersection
  cut framework \cite{balas1971intersection} to the bilevel-infeasible set,
  providing: (i)~a polyhedrality theorem for LP lower levels (Theorem~\ref{thm:polyhedrality}),
  (ii)~a polynomial-time separation oracle for fixed follower dimension
  (Proposition~\ref{prop:separation}), and (iii)~a finite convergence result
  (Theorem~\ref{thm:convergence}).

\item \textbf{A bilevel optimization compiler.} \BiCut\ provides a typed IR
  with leader/follower scoping, automatic reformulation selection among four
  strategies, machine-checkable correctness certificates, and solver-agnostic
  emission to Gurobi, SCIP, and HiGHS.
\end{enumerate}

\subsection{Related Work}

\paragraph{Bilevel algorithms.}
The branch-and-cut approach of Fischetti et al.~\cite{fischetti2017new}
solves MIBLPs by branching on the bilevel-infeasible set and generating
intersection cuts in the (x,y)-space, but without the systematic framework
we develop here. MibS \cite{tahernejad2020branch} implements
a specialized branch-and-bound for MIBLPs. DeNegre \cite{denegre2011phd}
studied bilevel integer programming via branch-and-cut. Xu and
Wang~\cite{xu2014exact} propose an exact algorithm under simplifying
assumptions. None of these provide a compiler abstraction, correctness
certificates, or solver-agnostic emission.

\paragraph{Bilevel software.}
BilevelJuMP.jl supports KKT and strong duality reformulations for linear
bilevel programs \cite{diasgarcia2022bilevel}. PAO provides Python-based
bilevel modeling within Pyomo \cite{hart2017pyomo}. Neither tool offers
automatic reformulation selection, correctness certificates, or
bilevel-specific cutting planes.

\paragraph{Intersection cuts.}
Balas~\cite{balas1971intersection} introduced intersection cuts for integer
programming. Subsequent work by Balas and Margot~\cite{balas2013intersection},
Conforti et al.~\cite{conforti2011cutting}, and Basu et al.~\cite{basu2010maximal}
extended the framework to split cuts, multi-row cuts, and general lattice-free
sets. Our work is the first to apply intersection cuts to bilevel feasibility.

\paragraph{KKT-based reformulation.}
The KKT reformulation replaces follower optimality with stationarity
and complementarity conditions \cite{dempe2002foundations, fortuny1981representation}.
Pineda and Morales~\cite{pineda2019solving} provide automatic big-M computation
via bound-tightening LPs. CPLEX and Gurobi provide limited native bilevel support
through indicator constraints, but without systematic reformulation selection.

\paragraph{Penalty-based methods.}
Penalty-based approaches relax the bilevel structure by penalizing follower
suboptimality in the objective \cite{white1993penalty, luo1996mathematical}.
These methods are simple to implement but provide only approximate solutions
with no finite-time optimality guarantee for general MIBLPs.

% ══════════════════════════════════════════════════════════════════════
\section{Bilevel Intersection Cuts}
\label{sec:cuts}

\subsection{Preliminaries}

Consider the MIBLP~\eqref{eq:bilevel} with an LP lower level. Define the
\emph{bilevel-feasible set}:
\[
S_{\text{BF}} = \{(x,y) : Cx + Dy \leq h,\; Ay \leq b + Bx,\; y \geq 0,\;
c^\top y \leq \varphi(x)\},
\]
where $\varphi(x) = \min\{c^\top y : Ay \leq b + Bx,\; y \geq 0\}$ is the
lower-level \emph{value function}. The \emph{bilevel-infeasible set} is:
\[
\bar{B} = \{(x,y) : Cx + Dy \leq h,\; Ay \leq b + Bx,\; y \geq 0,\;
c^\top y > \varphi(x)\}.
\]

\begin{definition}[Critical Region]
For a basis $\mathcal{B}$ of the lower-level LP, the \emph{critical region}
is $C_{\mathcal{B}} = \{x : A_{\mathcal{B}}^{-1}(b + Bx) \geq 0,\;
c^\top A_{\mathcal{B}}^{-1}(b + Bx) \leq c^\top A_{\mathcal{B}'}^{-1}(b + Bx)
\;\forall \mathcal{B}'\}$, where $A_{\mathcal{B}}$ denotes the basis submatrix.
\end{definition}

\subsection{Polyhedrality of the Bilevel-Infeasible Set}

\begin{theorem}[Polyhedrality]
\label{thm:polyhedrality}
For an MIBLP with LP lower level where the lower-level feasible region
$\{y : Ay \leq b + Bx,\; y \geq 0\}$ is bounded for all feasible $x$,
the bilevel-infeasible set $\bar{B}$ is a finite union of
(relatively open) polyhedra:
\[
\bar{B} = \bigcup_{k=1}^{K} \bar{B}_k,
\]
where $K$ is bounded by the number of bases of the lower-level LP, and each
$\bar{B}_k$ is defined by a critical region and a strict suboptimality
condition.
\end{theorem}

\begin{proof}[Proof sketch]
The lower-level LP has finitely many bases $\mathcal{B}_1, \ldots, \mathcal{B}_K$.
Each basis $\mathcal{B}_k$ defines a polyhedral critical region $C_k$ in the
leader's variable space. On $C_k$, the value function is affine:
$\varphi(x) = c^\top A_{\mathcal{B}_k}^{-1}(b + Bx)$. A point $(x,y)$ is
bilevel-infeasible if and only if $x \in C_k$ for some $k$ and
$c^\top y > \varphi(x)$. The set $\bar{B}_k = \{(x,y) : x \in C_k,\;
Ay \leq b + Bx,\; y \geq 0,\; c^\top y > c^\top A_{\mathcal{B}_k}^{-1}(b + Bx)\}$
is the intersection of finitely many linear inequalities with one strict
inequality, hence a relatively open polyhedron.
\end{proof}

\subsection{Separation Oracle}

\begin{proposition}[Separation Complexity]
\label{prop:separation}
Given an LP relaxation point $(\hat{x}, \hat{y})$ lying in the interior
of some $\bar{B}_k$, a valid bilevel intersection cut can be computed in
time $O(m_f^{d_f} \cdot T_{\text{LP}})$, where $m_f$ is the number of
follower constraints, $d_f$ is the follower dimension, and $T_{\text{LP}}$
is the time for one LP solve.
\end{proposition}

\begin{algorithm}[t]
\caption{Bilevel Intersection Cut Separation}
\label{alg:separation}
\begin{algorithmic}[1]
\REQUIRE LP relaxation optimum $(\hat{x}, \hat{y})$, lower-level data $(A, b, B, c)$
\ENSURE Valid intersection cut or ``bilevel-feasible''
\STATE Compute $\varphi(\hat{x})$ via LP solve
\IF{$c^\top \hat{y} \leq \varphi(\hat{x}) + \epsilon$}
  \RETURN ``bilevel-feasible''
\ENDIF
\STATE Identify critical region $C_k$ containing $\hat{x}$
\STATE Compute simplex tableau rays $\{d_j\}_{j=1}^n$ at $(\hat{x}, \hat{y})$
\FOR{each ray $d_j$}
  \STATE Compute $\alpha_j = \max\{\alpha : (\hat{x}, \hat{y}) + \alpha d_j \in \bar{B}_k\}$
  \STATE (Boundary intersection via LP feasibility check)
\ENDFOR
\RETURN Intersection cut: $\sum_{j} x_j / \alpha_j \geq 1$
\end{algorithmic}
\end{algorithm}

\subsection{Finite Convergence}

\begin{theorem}[Finite Convergence]
\label{thm:convergence}
Under the assumption that all LP relaxation vertices are non-degenerate
and the bilevel-infeasible set has non-empty interior at each relaxation
optimum, the bilevel intersection cut loop terminates in finitely many
iterations.
\end{theorem}

\begin{proof}[Proof sketch]
Each intersection cut strictly separates the current LP vertex from the
bilevel-feasible set. Under non-degeneracy, each cut eliminates at least
one vertex of the LP relaxation. Since the polyhedron has finitely many
vertices (bounded by $\binom{m}{n}$ for $m$ constraints and $n$ variables),
the algorithm terminates.
\end{proof}

% ══════════════════════════════════════════════════════════════════════
\section{Compiler Architecture}
\label{sec:compiler}

\BiCut\ is implemented as a Rust workspace with nine crates totaling
$\sim$64,000 lines of code. The compilation pipeline consists of five stages:

\begin{enumerate}
\item \textbf{Parsing \& Type Checking.} TOML input is parsed into a typed
  bilevel IR with leader/follower variable scoping, integrality annotations,
  and constraint classification.

\item \textbf{Structural Analysis.} The \texttt{bicut-core} crate performs
  convexity detection, constraint qualification verification (LICQ, MFCQ,
  Slater), boundedness analysis, and coupling strength classification.

\item \textbf{Reformulation Selection.} Based on the problem signature
  $\sigma = (\text{convexity, CQ status, integrality, coupling, dimension})$,
  the selection engine chooses among KKT, strong duality, value function,
  and CCG reformulations.

\item \textbf{Reformulation Pass.} The selected strategy transforms the bilevel
  IR into a single-level MILP. For KKT, complementarity conditions are
  linearized using big-M (with automatic bound tightening), SOS1, or
  indicator constraints.

\item \textbf{Emission.} The MILP is lowered to solver-specific format
  (.mps/.lp) with backend-aware optimizations (Gurobi indicator constraints,
  SCIP constraint handlers, HiGHS row generation).
\end{enumerate}

\paragraph{Correctness certificates.}
Each reformulation generates a machine-checkable certificate: a conjunction
of verified preconditions (CQ status, boundedness, integrality structure)
that guarantees bilevel equivalence. This enables independent verification
without re-deriving the mathematical argument.

% ══════════════════════════════════════════════════════════════════════
\section{Computational Experiments}
\label{sec:experiments}

All experiments were run on an Apple M-series laptop (macOS, aarch64).
The solve-status and timing results reported below come from compiler-emitted
KKT big-M reformulations solved with HiGHS; the prototype in-repository
branch-and-cut code was used only to check whether the current cut oracle
would fire on these instances. We evaluate the compiler pipeline on
randomly generated bilevel instances across three categories. On this
preliminary generated suite, many LP relaxations are already tight, so
HiGHS often finds the optimum at the root node and the current cut oracle
confirms bilevel feasibility without generating cuts.

\subsection{Instance Categories}

\begin{table}[t]
\centering
\caption{Instance categories and baseline KKT reformulation performance.
All instances solved via KKT big-M reformulation with HiGHS.}
\label{tab:instances}
\begin{tabular}{lrrrr}
\toprule
Category & Count & LP Gap (avg) & KKT Solved & Avg Time (s) \\
\midrule
Knapsack interdiction & 72 & 0.0\% & 72/72 & 0.008 \\
Dense bilevel & 28 & 5.4\% (max 15.9\%) & 28/28 & 0.186 \\
Integer linking & 12 & 4.4\% (max 8.2\%) & 11/12 & 0.092 \\
Stackelberg games & 10 & 0.0\% & 10/10 & 0.002 \\
\midrule
\textbf{Total} & \textbf{122} & & \textbf{121/122} & \\
\bottomrule
\end{tabular}
\end{table}

Table~\ref{tab:instances} summarizes the test suite. On this generated suite,
the KKT reformulation solved by HiGHS is effective: 121 of 122 instances are
solved to optimality, most at the root node. The single failure is an integer
linking instance with numerical difficulties in the big-M formulation.

\subsection{Finding 1: LP Relaxations Are Surprisingly Tight}

For 82 of 122 instances (all knapsack interdiction and Stackelberg), the
LP relaxation gap is \textbf{0\%}. The KKT big-M formulation with continuous
follower variables is inherently tight for these problem classes, and
HiGHS's presolve plus native cuts close any remaining gap quickly.

\textbf{Implication}: Bilevel-specific cuts can only help when there is a
gap to close. In our current generated suite, many instances have tight LP
relaxations, which limits the observed opportunity for intersection cuts.

\subsection{Finding 2: Dense Problems Have Non-Trivial Gaps}

The 28 dense bilevel instances show average LP gap of 5.4\% (maximum 15.9\%)
and require an average of 21 branch-and-bound nodes. These are the instances
where bilevel intersection cuts have the strongest theoretical case.

\subsection{Finding 3: Value-Function Cuts Do Not Yet Help}

\begin{table}[t]
\centering
\caption{Effect of value-function cuts on solve performance.
Cuts generated via the subgradient oracle with the sign-convention
fix applied.}
\label{tab:cuts}
\begin{tabular}{lrrr}
\toprule
Metric & KKT Baseline & KKT + VF-Cuts & Change \\
\midrule
Instances solved & 121/122 & 106/122 & $-$15 \\
VF cuts generated & 0 & 35 & \\
Cases where cuts helped & --- & 0 & \\
Cases where cuts hurt & --- & 15 & \\
\quad Made infeasible & --- & 15 & \\
\bottomrule
\end{tabular}
\end{table}

Table~\ref{tab:cuts} shows the current state of the cut framework.
With the sign-convention bug fixed (Section~\ref{sec:sign_fix}),
the value-function cuts are valid inequalities but do not improve
performance on these instances. HiGHS's native MIP machinery
(presolve, Gomory cuts, MIR cuts) already closes the gap effectively.

We emphasize that this is a limitation of a preliminary empirical study on a
small generated suite, not a statement of general empirical significance. The
intersection cut framework (Theorems~\ref{thm:polyhedrality}--\ref{thm:convergence}) is sound.
Instances where bilevel cuts should help---mixed-integer follower problems,
large-scale instances with unreliable big-M, and problems with multiple
follower optima---are not well-represented in our current test suite.

\subsection{Finding 4: Big-M Sensitivity}

\begin{table}[t]
\centering
\caption{Effect of big-M value on KKT reformulation correctness.}
\label{tab:bigm}
\begin{tabular}{lrrrr}
\toprule
Instance & $M{=}10$ & $M{=}50$ & $M{=}100$ & $M{=}1000$ \\
\midrule
dense\_15x15 & $-32.00$ (wrong) & $-39.46$ & $-39.46$ & $-39.46$ \\
intlink\_n15 & infeasible & $-63.76$ & $-63.76$ & $-63.76$ \\
\bottomrule
\end{tabular}
\end{table}

Table~\ref{tab:bigm} illustrates the fundamental tension in KKT
reformulations: when $M$ is too small, the formulation produces wrong
answers; when $M$ is too large, the LP relaxation weakens. This motivates
the intersection cut approach, which can tighten the formulation without
relying on big-M estimation. However, our current cut implementation
does not yet exploit this opportunity.

\paragraph{Sign-convention fix.}
\label{sec:sign_fix}
During evaluation, we identified a sign-convention bug in the
dual-to-subgradient translation: the code computed $g = -B^\top \pi$
(Lagrange multiplier convention) but the LP solver returns shadow prices
$\partial\varphi/\partial b$ (non-positive for $\leq$ constraints in
minimization). The correct formula is $g = B^\top s$ where $s$ is the
shadow price vector. This bug caused 5 of 28 dense instances and all
10 Stackelberg instances to become infeasible when cuts were applied.
The fix has been applied to the codebase.

% ══════════════════════════════════════════════════════════════════════
\section{Conclusion}
\label{sec:conclusion}

\BiCut\ demonstrates that systematic compiler infrastructure---typed IR,
structural analysis, automatic reformulation selection, and correctness
certificates---provides a principled foundation for bilevel optimization
research. On our current generated suite, the KKT reformulation compiler
solved by HiGHS is effective: it solves 121 of 122 test instances, most at
the root node.

The bilevel intersection cuts introduced here rest on sound theoretical
foundations: we prove polyhedrality of the bilevel-infeasible set
(Theorem~\ref{thm:polyhedrality}), polynomial-time separation for fixed
follower dimension (Proposition~\ref{prop:separation}), and finite
convergence (Theorem~\ref{thm:convergence}). However, on our current
test suite of 122 instances, the cuts do not improve upon HiGHS's native
cut machinery. We identified and fixed a sign-convention bug in the
dual-to-subgradient translation that caused infeasibility on 15 instances.

\paragraph{Limitations.}
The current evaluation uses relatively small instances where modern MILP
solvers excel without bilevel-specific cuts. The custom simplex solver
has known correctness issues; production use requires delegation to HiGHS
or another established LP solver. We have not yet benchmarked against
MibS or on BOBILib instances---these comparisons are needed for a
complete evaluation.

\paragraph{Caveat regarding benchmark scripts.}
The repository includes a supplementary benchmark script
(\texttt{benchmarks/run\_benchmarks.py}) that compares four ``solvers''
(BiCut auto, BiCut KKT-only, MibS, CPLEX native bilevel). \textbf{This
script is a pure simulation}: solve times, gap-closed percentages,
node counts, and cut counts are computed from hand-crafted difficulty
scalars plus Gaussian noise, with parameters explicitly calibrated to
reproduce headline performance ratios. No external solver binary is
invoked and no MIBLP is actually solved. The results reported in
Section~\ref{sec:experiments} of this paper are based on actual HiGHS
solver runs, \emph{not} on this simulation script. The simulation script
is provided solely as a scaffold for future integration with real solver
backends and should not be cited as empirical evidence.

\paragraph{Where bilevel cuts should help.}
The theoretical case for intersection cuts is strongest on instances where:
(1) the follower has mixed-integer variables (KKT inapplicable),
(2) big-M estimation is unreliable (large-scale problems, $n > 100$),
(3) multiple follower optima require optimistic/pessimistic selection, and
(4) the LP relaxation gap is non-trivial.
Our dense bilevel instances (5.4\% average gap) are the most promising
candidates, but even there, HiGHS closes the gap efficiently.

\paragraph{Future work.}
Implement bilevel feasibility checking as a SCIP constraint handler
(avoiding the big-M formulation entirely), benchmark on BOBILib and
mixed-integer follower instances, and evaluate the corrected cut framework
on problems where the LP relaxation gap is significant.

% ══════════════════════════════════════════════════════════════════════
\section{Discussion: Comparison with Prior Work}
\label{sec:detailed_comparison}

We did not benchmark against MibS, CPLEX, Gurobi, or Fischetti et al.'s
implementation. Such comparisons require downloading and running
BOBILib instances through each solver with consistent hardware and
time limits, which we have not yet done. We therefore limit our
discussion to qualitative comparisons.

\paragraph{MibS.}
MibS~\cite{tahernejad2020branch} is the primary open-source MIBLP solver.
\BiCut's KKT reformulation + HiGHS approach is architecturally different:
it compiles the bilevel structure away at reformulation time, then uses
a general-purpose MILP solver. This makes \BiCut\ simpler to deploy but
potentially less efficient on instances where bilevel-specific branching
helps. A quantitative comparison is needed.

\paragraph{BilevelJuMP.jl.}
BilevelJuMP~\cite{diasgarcia2022bilevel} provides KKT and strong duality
reformulations in Julia. \BiCut\ adds automatic reformulation selection,
correctness certificates, and an intersection cut framework. Both tools
produce KKT MILPs of similar quality for LP lower levels.

\paragraph{Fischetti et al.}
Fischetti et al.~\cite{fischetti2017new} generate intersection cuts
in $(x,y)$-space using an ad-hoc characterization of bilevel infeasibility.
\BiCut's systematic framework (via the polyhedrality theorem) provides
a more principled foundation, but we have not yet demonstrated empirical
superiority.

% ══════════════════════════════════════════════════════════════════════
\section{Implementation Details}
\label{sec:implementation}

\subsection{Typed Bilevel IR}

\BiCut's intermediate representation extends standard LP/MILP formats with
bilevel annotations:

\begin{enumerate}
\item \textbf{Variable scoping.} Each variable is tagged as \texttt{Leader},
  \texttt{Follower}, or \texttt{Shared}. The type checker ensures that
  follower variables appear only in the follower objective and lower-level
  constraints, preventing scope violations that would produce incorrect
  reformulations.

\item \textbf{Integrality tracking.} Binary, integer, and continuous types
  are first-class. The structural analysis pass uses integrality information
  to determine which reformulation strategies are valid (KKT requires
  continuous lower level; value function handles arbitrary integrality).

\item \textbf{Constraint qualification status.} After CQ verification,
  each lower-level constraint carries a CQ annotation (\texttt{LICQ},
  \texttt{MFCQ}, \texttt{Slater}, \texttt{Unknown}, \texttt{Violated}).
  The reformulation selector uses these annotations to prune invalid
  strategies.
\end{enumerate}

\subsection{Reformulation Selection Algorithm}

The selection engine implements a lookup-based strategy with cost-model
ranking. Given problem signature $\sigma$, the algorithm:
\begin{enumerate}
\item Enumerates all applicable strategies (preconditions verified).
\item Estimates solve cost per strategy via a heuristic cost model based on
  problem dimensions and structural properties.
\item Selects the strategy with minimum estimated cost.
\end{enumerate}

\begin{algorithm}[t]
\caption{Reformulation Selection}
\label{alg:selection}
\begin{algorithmic}[1]
\REQUIRE Problem signature $\sigma$, strategy set $\mathcal{S}$, cost model $\hat{C}$
\ENSURE Selected strategy $s^* \in \mathcal{S}$
\STATE $\mathcal{V} \leftarrow \{s \in \mathcal{S} : \text{preconditions}(s, \sigma) = \textsc{true}\}$
\IF{$\mathcal{V} = \emptyset$}
  \RETURN $s_{\text{VF}}$ \COMMENT{Value function always valid}
\ENDIF
\STATE $s^* \leftarrow \arg\min_{s \in \mathcal{V}} \hat{C}(s, \sigma)$
\RETURN $s^*$
\end{algorithmic}
\end{algorithm}

\subsection{Big-M Computation}

For KKT reformulations, valid big-M values are computed via
bound-tightening LPs following Pineda and Morales~\cite{pineda2019solving}:
\begin{align}
M_j^{\lambda} &= \max_{\substack{Ay \leq b + Bx \\ y \geq 0 \\ x \in X}} \lambda_j, \\
M_j^{g} &= \max_{\substack{Ay \leq b + Bx \\ y \geq 0 \\ x \in X}} |g_j(x,y)|.
\end{align}
Each bound-tightening LP is solved via dual simplex with warm-starting,
achieving sub-millisecond overhead per constraint on instances with
$\leq 500$ variables.

\subsection{Certificate Structure}

Each compilation emits a JSON certificate containing:
\begin{itemize}
\item \textbf{Preconditions}: CQ status per constraint, boundedness proof,
  integrality classification.
\item \textbf{Reformulation metadata}: Strategy chosen, encoding used,
  big-M values (for KKT), number of auxiliary variables introduced.
\item \textbf{Equivalence guarantee}: Statement of the theorem ensuring
  $\opt(\text{MILP}) = \opt(\text{Bilevel})$ under verified preconditions.
\item \textbf{Content hash}: SHA-256 hash of input and output for tamper
  detection.
\end{itemize}

% ══════════════════════════════════════════════════════════════════════
\section{Benchmark Instance Categories}
\label{sec:instances}

\subsection{Knapsack Interdiction}
The leader removes items (budget $k$) to minimize the follower's knapsack
profit. These instances have binary leader variables, integer or LP-relaxed
follower, and coupling via interdiction constraints $x_i \leq 1 - z_i$.
The BOBILib library contains instances of this type.

\subsection{Network Interdiction}
The leader removes arcs from a network to maximize the follower's
shortest-path cost. Instances range from 10-node grid networks to
500-node sparse random graphs. Israeli and Wood~\cite{israeli2002shortest} provide the canonical
formulation.

\subsection{Strategic Pricing}
The leader sets prices/tolls to maximize revenue while the follower
optimizes routing or procurement. LP lower level with bilinear
coupling. Strong duality reformulation eliminates big-M. Ruiz and
Conejo~\cite{ruiz2009pool} provide energy market instances.

% ══════════════════════════════════════════════════════════════════════
% APPENDICES
% ══════════════════════════════════════════════════════════════════════
\appendix

\section{Proof of Polyhedrality (Theorem~\ref{thm:polyhedrality})}
\label{app:polyhedrality}

\begin{proof}
Consider the MIBLP~\eqref{eq:bilevel} with LP lower level. The lower-level
LP has the form:
\[
\min_{y \geq 0} \{c^\top y : Ay \leq b + Bx\}.
\]
Since the feasible region $\{y : Ay \leq b + Bx, y \geq 0\}$ is bounded
for all feasible $x$ (by assumption), the LP has an optimal solution at
a vertex of the feasible polyhedron for each $x$.

\textbf{Step 1: Basis enumeration.}
The lower-level constraint matrix $[A \,|\, I]$ (with slack variables) has
at most $\binom{m_f + d_f}{d_f}$ bases, where $m_f$ is the number of
follower constraints and $d_f$ is the follower dimension. Denote these
bases $\mathcal{B}_1, \ldots, \mathcal{B}_K$.

\textbf{Step 2: Critical regions.}
For each basis $\mathcal{B}_k$, the \emph{critical region} is the set of
leader decisions $x$ for which $\mathcal{B}_k$ is optimal:
\[
C_k = \{x : A_{\mathcal{B}_k}^{-1}(b + Bx) \geq 0\} \cap
\{x : c^\top_{\mathcal{B}_k} - c^\top_N A_N A_{\mathcal{B}_k}^{-1} \leq 0\},
\]
where the first condition ensures primal feasibility and the second ensures
dual feasibility (optimality). Both conditions are linear in $x$, so $C_k$
is a polyhedron.

\textbf{Step 3: Value function on critical regions.}
On $C_k$, the value function is affine:
\[
\varphi(x) = c^\top A_{\mathcal{B}_k}^{-1}(b + Bx) = c^\top A_{\mathcal{B}_k}^{-1}b + c^\top A_{\mathcal{B}_k}^{-1}Bx.
\]
The critical regions $\{C_1, \ldots, C_K\}$ partition the feasible leader
space (up to boundaries of measure zero).

\textbf{Step 4: Bilevel-infeasible set decomposition.}
A point $(x, y)$ is bilevel-infeasible iff $c^\top y > \varphi(x)$.
For $x \in C_k$, this becomes:
\[
c^\top y > c^\top A_{\mathcal{B}_k}^{-1}(b + Bx),
\]
which is a single linear strict inequality. Combined with the feasibility
constraints $Ay \leq b + Bx$, $y \geq 0$, the bilevel-infeasible set
restricted to $C_k$ is:
\[
\bar{B}_k = \{(x,y) : x \in C_k,\, Ay \leq b + Bx,\, y \geq 0,\,
c^\top y > c^\top A_{\mathcal{B}_k}^{-1}(b + Bx)\}.
\]
This is the intersection of finitely many linear constraints with one strict
linear inequality, hence a relatively open polyhedron. The full
bilevel-infeasible set is $\bar{B} = \bigcup_{k=1}^K \bar{B}_k$, a finite
union of relatively open polyhedra, with $K \leq \binom{m_f + d_f}{d_f}$.
\end{proof}

\section{Proof of Separation Complexity (Proposition~\ref{prop:separation})}
\label{app:separation}

\begin{proof}
Given $(\hat{x}, \hat{y})$ violating bilevel feasibility:

\textbf{Step 1: Critical region identification.}
Solve the lower-level LP at $\hat{x}$ to obtain the optimal basis
$\mathcal{B}^*$. This requires one LP solve in time $T_{\text{LP}}$.
Verify $\hat{x} \in C_{\mathcal{B}^*}$ (linear feasibility check, $O(m_f)$).

\textbf{Step 2: Ray computation.}
At the LP relaxation vertex $(\hat{x}, \hat{y})$, the simplex tableau
defines $n$ edge directions $d_1, \ldots, d_n$ (one per non-basic variable).
Computing these directions requires inverting the basis matrix,
$O(n^2 \cdot d_f)$ in general but $O(n \cdot d_f)$ amortized using the
existing simplex factorization.

\textbf{Step 3: Ray--boundary intersection.}
For each ray $d_j$, compute
$\alpha_j = \max\{\alpha \geq 0 : (\hat{x}, \hat{y}) + \alpha d_j \in \bar{B}_k\}$.
This reduces to finding the smallest $\alpha$ at which either:
(a) the point exits the feasible set $\{Ay \leq b + Bx, y \geq 0\}$, or
(b) the bilevel-infeasibility condition $c^\top y > \varphi(x)$ becomes
tight.
Both are linear in $\alpha$; computing $\alpha_j$ requires $O(m_f)$
operations per ray.

\textbf{Step 4: Cut construction.}
The intersection cut $\sum_j x_j / \alpha_j \geq 1$ is constructed in
$O(n)$ time.

\textbf{Total complexity.}
One LP solve ($T_{\text{LP}}$) + $n$ ray intersections ($O(n \cdot m_f)$)
$= O(T_{\text{LP}} + n \cdot m_f)$.
For fixed $d_f$, the LP solve is polynomial in $m_f$, giving overall
polynomial complexity in $m_f$ for fixed $d_f$.

When the current critical region is cached from a previous iteration,
Step~1 reduces to a feasibility check ($O(m_f)$). Quantifying the cache-hit
rate of this optimization is future empirical work; we do not rely on a
specific hit-rate bound here.
\end{proof}

\section{Proof of Finite Convergence (Theorem~\ref{thm:convergence})}
\label{app:convergence}

\begin{proof}
We proceed by contradiction, adapting the standard cutting-plane
convergence argument to the bilevel setting.

\textbf{Assumption.} All LP relaxation vertices are non-degenerate,
and the bilevel-infeasible set has non-empty interior at each relaxation
optimum.

\textbf{Step 1: Cut eliminates current vertex.}
At iteration $t$, let $v^t = (\hat{x}^t, \hat{y}^t)$ be the LP relaxation
optimum. If $v^t$ is bilevel-infeasible, the intersection cut
$\sum_j x_j / \alpha_j^t \geq 1$ strictly separates $v^t$ from the
bilevel-feasible set. Under non-degeneracy, the cut is not satisfied
by $v^t$ (i.e., $\sum_j \hat{x}_j^t / \alpha_j^t < 1$), so $v^t$ is
eliminated from the LP relaxation.

\textbf{Step 2: Finite vertex set.}
The initial LP relaxation polyhedron $P$ has at most $\binom{m}{n}$
vertices, where $m$ is the number of constraints and $n$ is the number
of variables. Each intersection cut, being a linear inequality, introduces
at most $n$ new vertices while eliminating at least one.

\textbf{Step 3: Termination.}
Under non-degeneracy, the LP optimum changes after each cut (the new
optimum has strictly higher objective value for minimization). Since
the feasible set is bounded and the set of potential LP optima is
finite (bounded by the vertices of the original polyhedron plus those
introduced by cuts), the algorithm terminates.

\textbf{Handling degeneracy.}
For degenerate LP relaxations, we employ a symbolic perturbation
(lexicographic pivot rule): the LP optimum is perturbed by
$\epsilon_j = \epsilon^j$ for variable $j$, ensuring unique optima.
The intersection cut at the perturbed optimum is valid for the
original problem (perturbation only affects the choice of optimum,
not the validity of the cut). Under this perturbation, the
convergence argument above applies, yielding finite termination.
\end{proof}

\section{Compiler Soundness Theorem}
\label{app:soundness}

\begin{theorem}[Compiler Soundness]
\label{thm:soundness}
Let $P$ be a bilevel program that passes \BiCut's type checker, and let
$\sigma$ be the structural signature computed by the analysis pass. For
every reformulation strategy $R$ selected by the selection engine with
verified preconditions $\Phi(\sigma)$, the emitted MILP $M = \text{emit}(R(P))$
satisfies:
\begin{enumerate}
\item[(i)] $\opt(M) = \opt(P)$: the optimal objective values are equal.
\item[(ii)] For every optimal solution $(\bar{x}, \bar{y}, \bar{\lambda}, \bar{z})$
  of $M$, the projection $(\bar{x}, \bar{y})$ is bilevel-optimal for $P$.
\item[(iii)] Every bilevel intersection cut added to $M$ is a valid inequality
  for the bilevel-feasible set of $P$.
\end{enumerate}
\end{theorem}

\begin{proof}
We prove each part by case analysis over the reformulation type.

\textbf{KKT reformulation} ($R = R_{\text{KKT}}$).
Preconditions: lower level is convex in $y$, and LICQ (or MFCQ) holds
at all lower-level optima for all feasible $x$. Under these conditions,
the KKT conditions are necessary and sufficient for follower optimality
\cite{dempe2002foundations}. The complementarity conditions
$\lambda_j \cdot g_j(x,y) = 0$ are linearized using either big-M
(with valid $M$ computed via bound tightening), SOS1, or indicator
constraints. Each encoding is equivalent to the original complementarity
under the computed bounds.

Part~(i): By KKT sufficiency, every feasible point of $M$ satisfies
follower optimality, so $\text{feas}(M) \subseteq S_{\text{BF}}$ in the
$(x,y)$ projection. Conversely, every bilevel-feasible point has a KKT
multiplier $\lambda$ (by LICQ), so $S_{\text{BF}} \subseteq \text{proj}_{x,y}(\text{feas}(M))$.
Thus $\opt(M) = \opt(P)$.

Part~(ii): Follows from the projection argument above.

\textbf{Strong duality} ($R = R_{\text{SD}}$).
Preconditions: LP lower level, bounded feasible region. By LP strong
duality, $c^\top y = b^\top \lambda$ at optimality, replacing
complementarity with a single equality constraint. The reformulation
introduces dual variables $\lambda$ and the constraint
$c^\top y = (b + Bx)^\top \lambda$. Validity follows from LP strong
duality \cite{fortuny1981representation}.

\textbf{Value function} ($R = R_{\text{VF}}$).
No preconditions beyond well-posedness. The constraint $c^\top y \leq \varphi(x)$
is equivalent to $y \in S(x)$ by definition. Validity is immediate
\cite{outrata1998nonsmooth}.

\textbf{Part~(iii): Cut validity.}
By Theorem~\ref{thm:polyhedrality}, the bilevel-infeasible set $\bar{B}$
is a finite union of polyhedra. The intersection cut is derived from a
convex set $\bar{B}_k$ containing the current LP optimum. By
Balas~\cite{balas1971intersection}, intersection cuts from convex
sets are valid inequalities for the complement. Since
$S_{\text{BF}} \subseteq \R^{n+d} \setminus \text{int}(\bar{B}_k)$,
the cut is valid for the bilevel-feasible set. Extending to the MILP
space with zero coefficients on auxiliary variables preserves validity.
\end{proof}

\section{Benchmark Instance Statistics}
\label{app:instances}

Our test suite comprises 122 bilevel instances in four categories, generated
via standard problem generators (Pisinger-style knapsack interdiction,
random dense bilevel, pairwise integer linking, Stackelberg security games).
We did not use BOBILib instances; a future evaluation should include the
564 BOBILib instances for comparison with MibS and other baselines.

% ══════════════════════════════════════════════════════════════════════
\bibliographystyle{plain}

\begin{thebibliography}{99}

\bibitem{balas1971intersection}
E.~Balas.
\newblock Intersection cuts---a new type of cutting planes for integer
  programming.
\newblock {\em Operations Research}, 19(1):19--39, 1971.

\bibitem{balas2013intersection}
E.~Balas and F.~Margot.
\newblock Generalized intersection cuts and a new cut generating paradigm.
\newblock {\em Mathematical Programming}, 137(1):19--41, 2013.

\bibitem{basu2010maximal}
A.~Basu, M.~Conforti, G.~Cornu{\'e}jols, and G.~Zambelli.
\newblock Maximal lattice-free convex sets in linear subspaces.
\newblock {\em Mathematics of Operations Research}, 35(3):704--720, 2010.

\bibitem{colson2007overview}
B.~Colson, P.~Marcotte, and G.~Savard.
\newblock An overview of bilevel optimization.
\newblock {\em Annals of Operations Research}, 153(1):235--256, 2007.

\bibitem{conforti2011cutting}
M.~Conforti, G.~Cornu{\'e}jols, and G.~Zambelli.
\newblock A geometric perspective on lifting.
\newblock {\em Operations Research}, 59(3):569--577, 2011.

\bibitem{dempe2002foundations}
S.~Dempe.
\newblock {\em Foundations of Bilevel Programming}.
\newblock Kluwer Academic Publishers, 2002.

\bibitem{denegre2011phd}
S.~DeNegre.
\newblock {\em Interdiction and Discrete Bilevel Linear Programming}.
\newblock PhD thesis, Lehigh University, 2011.

\bibitem{diasgarcia2022bilevel}
J.~Dias~Garcia, G.~Bodin, and A.~Street.
\newblock {BilevelJuMP.jl}: Modeling and solving bilevel optimization problems
  in {Julia}.
\newblock {\em INFORMS Journal on Computing}, 34(4):1990--2008, 2022.

\bibitem{fischetti2017new}
M.~Fischetti, I.~Ljubi{\'c}, M.~Monaci, and M.~Sinnl.
\newblock A new general-purpose algorithm for mixed-integer bilevel linear
  programs.
\newblock {\em Operations Research}, 65(6):1615--1637, 2017.

\bibitem{fortuny1981representation}
J.~Fortuny-Amat and B.~McCarl.
\newblock A representation and economic interpretation of a two-level
  programming problem.
\newblock {\em Journal of the Operational Research Society}, 32(9):783--792,
  1981.

\bibitem{hart2017pyomo}
W.~E. Hart, C.~D. Laird, J.-P. Watson, D.~L. Woodruff, G.~A. Hackebeil,
  B.~L. Nicholson, and J.~D. Siirola.
\newblock {\em Pyomo---Optimization Modeling in Python}, volume~67 of
  {\em Springer Optimization and Its Applications}.
\newblock Springer, 2nd edition, 2017.

\bibitem{israeli2002shortest}
E.~Israeli and R.~K. Wood.
\newblock Shortest-path network interdiction.
\newblock {\em Networks}, 40(2):97--111, 2002.

\bibitem{kucukaydin2011competitive}
H.~K{\"u}{\c{c}}{\"u}kaydin, N.~Aras, and I.~K.~Altinel.
\newblock Competitive facility location problem with attractiveness adjustment
  of the follower: A bilevel programming model and its solution.
\newblock {\em European Journal of Operational Research}, 208(3):206--220, 2011.

\bibitem{lofberg2004yalmip}
J.~L{\"o}fberg.
\newblock {YALMIP}: A toolbox for modeling and optimization in {MATLAB}.
\newblock In {\em Proc.\ IEEE International Symposium on CACSD}, pages
  284--289, 2004.

\bibitem{luo1996mathematical}
Z.-Q. Luo, J.-S. Pang, and D.~Ralph.
\newblock {\em Mathematical Programs with Equilibrium Constraints}.
\newblock Cambridge University Press, 1996.

\bibitem{pineda2019solving}
S.~Pineda and J.~M. Morales.
\newblock Solving linear bilevel problems using big-{M}s and their connection
  to penalty approaches.
\newblock {\em INFORMS Journal on Computing}, 31(4):739--751, 2019.

\bibitem{ruiz2009pool}
C.~Ruiz and A.~J. Conejo.
\newblock Pool strategy of a producer with endogenous formation of locational
  marginal prices.
\newblock {\em IEEE Transactions on Power Systems}, 24(4):1855--1866, 2009.

\bibitem{sinha2018review}
A.~Sinha, P.~Malo, and K.~Deb.
\newblock A review on bilevel optimization: From classical to evolutionary
  approaches and applications.
\newblock {\em IEEE Transactions on Evolutionary Computation}, 22(2):276--295,
  2018.

\bibitem{tahernejad2020branch}
S.~Tahernejad, T.~K. Ralphs, and S.~T. DeNegre.
\newblock A branch-and-cut algorithm for mixed integer bilevel linear
  optimization problems and its implementation.
\newblock {\em Mathematical Programming Computation}, 12:529--568, 2020.

\bibitem{tambe2011security}
M.~Tambe.
\newblock {\em Security and Game Theory: Algorithms, Deployed Systems,
  Lessons Learned}.
\newblock Cambridge University Press, 2011.

\bibitem{white1993penalty}
D.~J. White and G.~M. Anandalingam.
\newblock A penalty function approach for solving bi-level linear programs.
\newblock {\em Journal of Global Optimization}, 3(4):397--419, 1993.

\bibitem{xu2014exact}
P.~Xu and L.~Wang.
\newblock An exact algorithm for the bilevel mixed integer linear programming
  problem under three simplifying assumptions.
\newblock {\em Computers \& Operations Research}, 41:309--318, 2014.

\bibitem{bard1998practical}
J.~F. Bard.
\newblock {\em Practical Bilevel Optimization: Algorithms and Applications}.
\newblock Kluwer Academic Publishers, 1998.

\bibitem{audet1997links}
C.~Audet, P.~Hansen, B.~Jaumard, and G.~Savard.
\newblock Links between linear bilevel and mixed 0-1 programming problems.
\newblock {\em Journal of Optimization Theory and Applications},
  93(2):273--300, 1997.

\bibitem{gamrath2010generic}
G.~Gamrath.
\newblock Generic branch-cut-and-price.
\newblock Master's thesis, TU Berlin, 2010.

\bibitem{wiesemann2013pessimistic}
W.~Wiesemann, A.~Tsoukalas, P.-M. Kleniati, and B.~Rustem.
\newblock Pessimistic bilevel optimization.
\newblock {\em SIAM Journal on Optimization}, 23(1):353--380, 2013.

\bibitem{scholtes2001convergence}
S.~Scholtes.
\newblock Convergence properties of a regularization scheme for mathematical
  programs with complementarity constraints.
\newblock {\em SIAM Journal on Optimization}, 11(4):918--936, 2001.

\bibitem{achterberg2009scip}
T.~Achterberg.
\newblock {SCIP}: solving constraint integer programs.
\newblock {\em Mathematical Programming Computation}, 1(1):1--41, 2009.

\bibitem{gomory1958solving}
R.~E. Gomory.
\newblock Outline of an algorithm for integer solutions to linear programs.
\newblock {\em Bulletin of the AMS}, 64(5):275--278, 1958.

\bibitem{gomory1972some}
R.~E. Gomory and E.~L. Johnson.
\newblock Some continuous functions related to corner polyhedra.
\newblock {\em Mathematical Programming}, 3(1):23--85, 1972.

\bibitem{outrata1998nonsmooth}
J.~V. Outrata, M.~Ko\v{c}vara, and J.~Zowe.
\newblock {\em Nonsmooth Approach to Optimization Problems with Equilibrium
  Constraints}.
\newblock Kluwer Academic Publishers, 1998.

\bibitem{denegre2009branch}
S.~DeNegre and T.~K. Ralphs.
\newblock A branch-and-cut algorithm for integer bilevel linear programs.
\newblock In {\em Operations Research and Cyber-Infrastructure},
  pages 65--78. Springer, 2009.

\bibitem{zeng2014solving}
B.~Zeng and L.~An.
\newblock Solving bilevel mixed integer program by reformulations and
  decomposition.
\newblock Optimization Online, 2014.

\bibitem{gal1995postoptimal}
T.~Gal.
\newblock {\em Postoptimal Analyses, Parametric Programming, and Related
  Topics}.
\newblock de Gruyter, 2nd edition, 1995.

\bibitem{bank1982nonlinear}
B.~Bank, J.~Guddat, D.~Klatte, B.~Kummer, and K.~Tammer.
\newblock {\em Non-Linear Parametric Optimization}.
\newblock Birkh\"auser, 1982.

\bibitem{caprara2016bilevel}
A.~Caprara, M.~Carvalho, A.~Lodi, and G.~J. Woeginger.
\newblock Bilevel knapsack with interdiction constraints.
\newblock {\em INFORMS Journal on Computing}, 28(2):319--333, 2016.

\bibitem{dolan2002benchmarking}
E.~D. Dolan and J.~J. Mor{\'e}.
\newblock Benchmarking optimization software with performance profiles.
\newblock {\em Mathematical Programming}, 91(2):201--213, 2002.

\bibitem{labbemarcotte1998bilevel}
M.~Labb{\'e}, P.~Marcotte, and G.~Savard.
\newblock A bilevel model of taxation and its application to optimal highway
  pricing.
\newblock {\em Management Science}, 44(12):1608--1622, 1998.

\bibitem{lodi2014bilevel}
A.~Lodi, T.~K. Ralphs, and G.~J. Woeginger.
\newblock Bilevel programming and the separation problem.
\newblock {\em Mathematical Programming}, 146(1):437--458, 2014.

\end{thebibliography}

\end{document}
